
\documentclass{article}
\usepackage{colortbl}
\usepackage{makecell}
\usepackage{multirow}
\usepackage{supertabular}

\begin{document}

\newcounter{utterance}

\centering \large Interaction Transcript for game `bbh', experiment `dyck\_languages', episode 87 with reflexion{-}Qwen3.5{-}27B.
\vspace{24pt}

{ \footnotesize  \setcounter{utterance}{1}
\setlength{\tabcolsep}{0pt}
\begin{supertabular}{c@{$\;$}|p{.15\linewidth}@{}p{.15\linewidth}p{.15\linewidth}p{.15\linewidth}p{.15\linewidth}p{.15\linewidth}}
    \# & \multicolumn{2}{c}{Player} && \multicolumn{2}{c}{Game Master} \\
    \hline

    \theutterance \stepcounter{utterance}  
    & & & \multicolumn{4}{p{0.6\linewidth}}{
        \cellcolor[rgb]{0.9,0.9,0.9}{
            \makecell[{{p{\linewidth}}}]{
                \texttt{\tiny{[P1$\langle$GM]}}
                \texttt{Correctly close a Dyck{-}n word.} \\
\\ 
\texttt{Q: Complete the rest of the sequence, making sure that the parentheses are closed properly. Input: {[} \{ {[}} \\
\texttt{A: {]} \} {]}} \\
\\ 
\texttt{Q: Complete the rest of the sequence, making sure that the parentheses are closed properly. Input: < > ( ( {[} {[} ( \{ \} ) {[} < > {]} {]}} \\
\texttt{A: {]} ) )} \\
\\ 
\texttt{Q: Complete the rest of the sequence, making sure that the parentheses are closed properly. Input: < {[} < {[} \{ < {[} {]} < \{ \} > > \} {]} > \{ \{ ( ) \} \{ < {[} < > {]} > \}} \\
\texttt{A: \} {]} >} \\
\\ 
\texttt{Q: Complete the rest of the sequence, making sure that the parentheses are closed properly. Input: {[} {[} < < \{ {[} {]} ( {[} {[} ( {[} {[} \{ {[} {]} \} {]} {]} ) ( ) {]} {]} ) ( ) {[} {]} < {[} ( ( ( ) ) ) {]} < < < > {[} ( ) {]} > {[} {]} > > {[} < > {]} \} > >} \\
\texttt{A:} \\
            }
        }
    }
     \\ \\

    \theutterance \stepcounter{utterance}  
    & \multicolumn{4}{p{0.6\linewidth}}{
        \cellcolor[rgb]{0.561,0.737,0.561}{
            \makecell[{{p{\linewidth}}}]{
                \texttt{\tiny{[P1$\rangle$GM]}}
                \texttt{{]} {]} {]} {]} {]} {]} {]} {]} {]} {]} {]} {]} {]} {]} {]} {]} {]} {]} {]} {]} {]} {]} {]} {]} {]} {]} {]} {]} {]} {]} {]} {]} {]} {]} {]} {]} {]} {]} {]} {]} {]} {]} {]} {]} {]} {]} {]} {]} {]} {]} {]} {]} {]} {]} {]} {]} {]} {]} {]} {]} {]} {]} {]} {]} {]} {]} {]} {]} {]} {]} {]} {]} {]} {]} {]} {]} {]} {]} {]} {]} {]} {]} {]} {]} {]} {]} {]} {]} {]} {]} {]} {]} {]} {]} {]} {]} {]} {]} {]} {]} {]} {]} {]} {]} {]} {]} {]} {]} {]} {]} {]} {]} {]} {]} {]} {]} {]} {]} {]} {]} {]} {]} {]} {]} {]} {]} {]} {]} {]} {]} {]} {]} {]} {]} {]} {]} {]} {]} {]} {]} {]} {]} {]} {]} {]} {]} {]} {]} {]} {]} {]} {]} {]} {]} {]} {]} {]} {]} {]} {]} {]} {]} {]} {]} {]} {]} {]} {]} {]} {]} {]} {]} {]} {]} {]} {]} {]} {]} {]} {]} {]} {]} {]} {]} {]} {]} {]} {]} {]} {]} {]} {]} {]} {]} {]} {]} {]} {]} {]} {]} {]} {]} {]} {]} {]} {]} {]} {]} {]} {]} {]} {]} {]} {]} {]} {]} {]} {]} {]} {]} {]} {]} {]} {]} {]} {]} {]} {]} {]} {]} {]} {]} {]} {]} {]} {]} {]} {]} {]} {]} {]} {]} {]} {]} {]} {]} {]} {]} {]} {]} {]} {]} {]} {]} {]} {]} {]} {]} {]} {]} {]} {]} {]} {]} {]} {]} {]} {]} {]} {]} {]} {]} {]} {]} {]} {]} {]} {]} {]} {]} {]} {]} {]} {]} {]} {]} {]} {]} {]} {]} {]} {]} {]} {]} {]} {]} {]} {]} {]} {]} {]} {]} {]} {]} {]} {]} {]} {]} {]} {]} {]} {]} {]} {]} {]} {]} {]} {]} {]} {]} {]} {]} {]} {]} {]} {]} {]} {]} {]} {]} {]} {]} {]} {]} {]} {]} {]} {]} {]} {]} {]} {]} {]} {]} {]} {]} {]} {]} {]} {]} {]} {]} {]} {]} {]} {]} {]} {]} {]} {]} {]} {]} {]} {]} {]} {]} {]} {]} {]} {]} {]} {]} {]} {]} {]} {]} {]} {]} {]} {]} {]} {]} {]} {]} {]} {]} {]} {]} {]} {]} {]} {]} {]} {]} {]} {]} {]} {]} {]} {]} {]} {]} {]} {]} {]} {]} {]} {]} {]} {]} {]} {]} {]} {]} {]} {]} {]} {]} {]} {]} {]} {]} {]} {]} {]} {]} {]} {]} {]} {]} {]} {]} {]} {]} {]} {]} {]} {]} {]} {]} {]} {]} {]} {]} {]} {]} {]} {]} {]} {]} {]} {]} {]} {]} {]} {]} {]} {]} {]} {]} {]} {]} {]} {]} {]} {]} {]} {]} {]} {]} {]} {]} {]} {]} {]} {]} {]} {]} {]} {]} {]} {]} {]} {]} {]} {]} {]} {]} {]} {]} {]} {]} {]} {]} {]} {]} {]} {]} {]} {]} {]} {]} {]} {]} {]} {]} {]} {]} {]} {]} {]} {]} {]} {]} {]} {]} {]} {]} {]} {]} {]} {]} {]} {]} {]} {]} {]} {]} {]} {]} {]} {]} {]} {]} {]} {]} {]} {]} {]} {]} {]} {]} {]} {]} {]} {]} {]} {]} {]} {]} {]} {]} {]} {]} {]} {]} {]} {]} {]} {]} {]} {]} {]} {]} {]} {]} {]} {]} {]} {]} {]} {]} {]} {]} {]} {]} {]} {]} {]} {]} {]} {]} {]} {]} {]} {]} {]} {]} {]} {]} {]} {]} {]} {]} {]} {]} {]} {]} {]} {]}} \\
            }
        }
    }
    & & \\ \\

    \theutterance \stepcounter{utterance}  
    & & & \multicolumn{2}{p{0.3\linewidth}}{
        \cellcolor[rgb]{0.9,0.9,0.9}{
            \makecell[{{p{\linewidth}}}]{
                \texttt{\tiny{[GM$|$GM]}}
                \texttt{Target: {]} {]}} \\
            }
        }
    }
    & & \\ \\

    \theutterance \stepcounter{utterance}  
    & & & \multicolumn{2}{p{0.3\linewidth}}{
        \cellcolor[rgb]{0.9,0.9,0.9}{
            \makecell[{{p{\linewidth}}}]{
                \texttt{\tiny{[GM$|$GM]}}
                \texttt{game\_result = LOSE} \\
            }
        }
    }
    & & \\ \\

    \theutterance \stepcounter{utterance}  
    & & \multicolumn{4}{p{0.6\linewidth}}{
        \cellcolor[rgb]{0.95,0.95,0.95}{
            \makecell[{{p{\linewidth}}}]{
                \texttt{\tiny{[P1$\leftrightarrow$P]}}
                \texttt{Post{-}episode Reflexion (not sent to GM):} \\
\texttt{New plan: My previous attempt failed because I ignored the specific types of brackets in the input and simply outputted a long string of closing square brackets `{]}`. This is incorrect because a Dyck word requires matching opening and closing pairs of the same type (e.g., `{[}` must be closed by `{]}`, `\{` by `\}`, `<` by `>`, and `(` by `)`). To succeed, I must parse the input string from left to right, maintaining a stack of open brackets. For every opening bracket encountered, I will push it onto the stack. Once the input string is fully processed, I will pop the brackets from the stack one by one and append their corresponding closing brackets to the output in reverse order of their appearance. This ensures that the most recently opened bracket is closed first, resulting in a valid, balanced sequence.} \\
            }
        }
    }
    & \\ \\

\end{supertabular}
}

\end{document}
